
\documentclass[oneside]{ausarbeitung}
\bibliography{latexlit}


\emergencystretch 3em



% ----------------------------------------------------------------------

\begin{document}
	
	%--- Sprachauswahl
	% Erlaubte Werte:
	%   \selectlanguage{english}
	%   \selectlanguage{ngerman}
	\selectlanguage{ngerman}
	
	%--- Art der Arbeit
	% Erlaubte Werte:
	%   \Praxissemesterbericht
	%   \Projektbericht
	%   \Bachelorarbeit
	%   \Seminararbeit
	   \Masterarbeit
	
	%\Praxissemesterbericht
	
	%--- Studiengang:
	% Erlaubte Werte:
	%   \Informatik
	%   \Elektronik
	%   \DataScience
	\Informatik
	
	\title{JSONPath-Schnittstelle für NoSQL-Datenbanken}
	
	\author{Paul Mayr}
	\matrikelnr{77868}
	
	%--- Ist der Erstbetreuer (\examinerA) an der Hochschule ein Professor?
	% Erlaubte Werte:
	\examinerIsAProfessortrue   % Ja
	%   \examinerIsAProfessorfalse  % Nein
	%\examinerIsAProfessortrue   % Ja
	
	%--- Betreuer
	\examinerA{Prof.~Dr.~Winfried~Bantel}
	\examinerB{Prof.~Dr.~habil.~Christian~Heinlein}
	
	%--- Einreichungsdatum
	\date{01. Juni 2026}
	
	%--- Angaben zur Firma
	% Auskommentieren, wenn die Arbeit nicht bei einer ext. Firma gemacht wurde.
	
	%--- Titelseite Anzeigen
	\maketitle
	\cleardoublepage
	
	%---
	\pagenumbering{roman}
	\setcounter{page}{1}
	
	%--- Firmendaten Anzeigen
	% Auskommentieren, wenn die Arbeit nicht bei einer ext. Firma gemacht wurde.
	
	
	
	%---
	
	%-----------------------------------------------------------------------
	\cleardoublepage
	\tableofcontents
	
	
	
	
	\cleardoublepage
	\pagenumbering{arabic}
	\setcounter{page}{1}
	
	% ----------------------------------------------------------------------




	\chapter{Einleitung}
	
	\section{Motivation}
	
	\section{Problemstellung und Abgrenzung}
	
	\section{Ziel der Arbeit}
	
	\section{Vorgehen}
	
	
	
	\chapter{Grundlagen}
	
	\section{JSONPath}
	
	
	
	\chapter{Problemanalyse}
	
	
	
	\chapter{Lösungskonzept}
	
	
	
	\chapter{Implementierung}
	
	
	
	\chapter{Evaluierung}
	
	
	
	\chapter{Zusammenfassung und Ausblick}
	
	\section{Erreichte Ergebnisse}
	
	\section{Ausblick}
	
	%----------------------------------------------------------------------
	
	%---
	%\chapter{Anhang A}
	
	%---
	%\chapter{Anhang B}
	
	
\end{document}